\section{Collision invariants and closed dissipation forms for one-dimensional square-root dispersion}
\label{pin:section}

We consider the four-wave resonance relation associated with a pinned
one-dimensional chain. The angular variable and the usual lattice momentum
are related by $x=2\pi k$. Fix
\begin{equation}
 \mathbb T_{2\pi}=\R/(2\pi\mathbb Z),\qquad
 0<d<\frac12,\qquad
 \omega(x)=\sqrt{1-2d\cos x},\qquad
 m(\dd x)=\frac{\dd x}{2\pi}.
 \label{pin:dispersion}
\end{equation}
The dispersion parameter $d$ is unrelated to the small parameter in the
macroscopic limit considered later. With $d$ fixed, we write
$\omega=\omega_d$, restoring the subscript when comparing parameters.
This dispersion and its four-wave collision operator appear in
\cite{ALS2006,Luk2016}; the collision invariant problem is also discussed
in \cite{GLM2026}. Every result below holds for each fixed
$0<d<1/2$, with constants that may depend on $d$.

A collision consists of two incoming and two outgoing momenta, which we
also call its four legs. The invariant relation compares the values of
one function at these four momenta. We use two features of this relation.
On a nondegenerate resonance chart, the value at the target point equals
an average of the signed sum of the other three values. A change of
integration variable transfers derivatives in the target variable to a
smooth integral kernel. This yields regularity of invariants and
compactness for approximate invariants. On a critical chart, the energy
difference and the collision difference share a double factor; its
cancellation gives finite dissipation for smooth functions. We proceed
through the construction of resonance charts, the classification of
invariants, the positive spectral gap, and the closed-form realization.

\subsection{Periodic resonance, coarea measure, and the classification theorem}
\label{pin:model}

The four momenta satisfy
\begin{equation}
 x_0+x_1=x_2+x_3\pmod{2\pi},\qquad
 \omega(x_0)+\omega(x_1)=\omega(x_2)+\omega(x_3).
 \label{pin:resonance}
\end{equation}
The first two positions $x_0,x_1$ are incoming and the last two
$x_2,x_3$ are outgoing. Periodic momentum conservation includes both
normal and Umklapp processes: for chosen real representatives, the two
momentum sums may differ by an integer multiple of $2\pi$.
Eliminate the fourth variable, use $(x,y,z)=(x_0,x_1,x_2)$ as coordinates,
and set
\begin{align}
 w&=x+y-z\pmod{2\pi},\notag\\
 F(x,y,z)&=\omega(x)+\omega(y)-\omega(z)-\omega(w),\notag\\
 \Sigma&=\{(x,y,z)\in\mathbb T_{2\pi}^3:F=0,\ \nabla F\ne0\},
 \qquad
 \dd\Xi=(2\pi)^{-3}\frac{\dd\mathcal H^2}{|\nabla F|}.
 \label{pin:measure}
\end{align}
Here $\mathcal H^2$ is Euclidean two-dimensional area in local real
coordinates, and $\nabla F$ is the gradient in the same coordinates.
The coarea factor $|\nabla F|^{-1}$ restricts the three-dimensional
momentum integral to the zero-energy surface; $(2\pi)^{-3}$ comes from
the probability measure normalization of the three remaining angular
variables. The measure $\Xi=\Xi_d$ includes all regular resonance
branches. It is locally finite and $\sigma$-finite on $\Sigma$, but its
total mass need not remain finite near critical points of the energy.
For a periodic function $\varphi$, define the collision difference by
\begin{equation}
 \Delta\varphi(x,y,z)
 =\varphi(x)+\varphi(y)-\varphi(z)-\varphi(x+y-z).
 \label{pin:difference}
\end{equation}
A resonance is trivial if the two pairs of momenta differ only by a
permutation; its collision difference vanishes identically. Retaining
all four terms serves two purposes: it annihilates both constants and
frequency on resonance, and it cancels their common lower-order parts
near critical points. The regularization and dissipation estimates
below concern this collision difference.

\begin{theorem}[Finite measurable collision invariants]
\label{pin:classification}
Let $\varphi:\mathbb T_{2\pi}\to\C$ be measurable and finite almost
everywhere, and suppose that
\begin{equation}
                    \Delta\varphi=0\qquad\Xi\text{-almost everywhere}.
 \label{pin:invariance}
\end{equation}
There are unique $A,B\in\C$ such that
\begin{equation}
                      \varphi=A+B\omega\qquad m\text{-almost everywhere}.
 \label{pin:classified}
\end{equation}
In particular, $\varphi$ has a unique smooth periodic representative.
Conversely, every affine-frequency representative
$x\mapsto A+B\omega(x)$ satisfies the invariant identity on every
resonance in \eqref{pin:resonance}. Real-valued invariants have real
coefficients $A,B$.
\end{theorem}

Regularity in this theorem is a consequence of the resonance identity.
For a measurable function that is merely finite almost everywhere, we
first compare level sets along the three non-target legs to obtain an
$L^\infty$ bound on the circle. Resonance averages then produce a
smooth representative. Only at that stage do we differentiate along
resonance charts: second-order compatibility for the odd part and
third-order compatibility for the even part, together with periodic
matching, determine $A,B$. We first verify that changing the function
on a one-dimensional null set leaves its identity on the resonance
surface unchanged.

\begin{lemma}[Absolute continuity of the four coordinate maps]
\label{pin:leg-ac}
Let $E\subset\mathbb T_{2\pi}$ be a one-dimensional Lebesgue null set.
For each of the maps $\pi_i=x,y,z,x+y-z$, one has
\begin{equation}
                         \Xi(\pi_i^{-1}(E))=0.
 \label{pin:ac}
\end{equation}
Thus \eqref{pin:invariance} is independent of the almost-everywhere
representative of $\varphi$.
\end{lemma}

\begin{proof}
The surface $\Sigma$ is a real-analytic two-dimensional manifold.
Choose a countable cover by relatively compact connected analytic
charts. We first show that $x$ is nonconstant on every nonempty open
surface patch. If $x=x_*$ on such a patch, its tangent space is
$\{\dd x=0\}$, so $(y,z)$ gives local coordinates on the patch.
Consequently $F(x_*,y,z)=0$ on a two-dimensional open set. Taking the
mixed derivative in $y,z$ gives
\[
                   \omega''(x_*+y-z)=0
\]
on a real interval. Analyticity forces $\omega''$ to vanish identically,
and periodicity then makes $\omega$ constant, contrary to $d>0$.

On each connected analytic chart, therefore, some coordinate derivative
of $x$ is not identically zero. The zero set of a nonzero real-analytic
function is Lebesgue null, so the set where $\dd x=0$ has zero
two-dimensional measure. The analytic fact used here follows from local
power series and Fubini: choose a nonzero coefficient in an expansion
in one variable; outside a null set of parameters, the resulting
one-variable analytic function has discrete zeros, and a countable
rectangle cover completes the argument. Away from this null set,
$x$ is a local submersion. Shrinking the chart and changing a coordinate
whose derivative is bounded away from zero proves \eqref{pin:ac} for
$\pi_0=x$.

The other three legs follow by exchanging the incoming legs, the
outgoing legs, or the incoming and outgoing pairs. In the three
coordinates left after elimination, these are integer linear
transformations with determinant of absolute value $1$. They preserve
$F$ up to sign, and hence preserve regularity. On every relatively
compact regular chart, $\Xi$ has a smooth positive density with respect
to two-dimensional coordinate measure. A countable cover gives the
claim on all of $\Sigma$.
\end{proof}

We will also prove a positive spectral gap for the collision difference,
finite energy of smooth functions, and a maximal closed form.
We begin with resonance coordinates covering every point of the circle.

\subsection{Nondegenerate resonance charts and regularization of measurable functions}
\label{pin:atlas}

Write $v=\omega'$ for the group velocity. We use the first leg $x$ as
the target variable and $z$ as the averaging variable, determining
$y,w$ from the conservation laws. Suppose a resonant quartet
$(x_*,y_*,z_*,w_*)$ has pairwise distinct non-target velocities
$v(y_*),v(z_*),v(w_*)$. Then
$F_y=v(y_*)-v(w_*)\ne0$, and the implicit function theorem gives a
resonance chart
\begin{equation}
 (x,H(x,z),z,W(x,z)),\qquad W=x+H-z,
 \label{pin:chart}
\end{equation}
on a real coordinate rectangle $I\times J$. Differentiating the
energy identity gives
\begin{equation}
 H_z=\frac{v(z)-v(W)}{v(H)-v(W)},\qquad
 W_z=\frac{v(z)-v(H)}{v(H)-v(W)}.
 \label{pin:chart-derivatives}
\end{equation}
The nonzero denominator allows us to solve the energy equation for
$y=H(x,z)$; the two nonzero numerators allow both $H,W$ to be used
as one-dimensional changes of integration variable. These three
nondegeneracy conditions have distinct roles. After shrinking the
rectangle, each leg takes values in a real lift interval of length
less than $2\pi$, and $H_z,W_z$ and $v(H)-v(W)$ are bounded away
from zero on the closure. For fixed $x$, the maps $H,W$ are therefore
monotone in $z$ and injective even as circle-valued maps. The coarea
measure on this chart is
\begin{equation}
              \dd\Xi=(2\pi)^{-3}
                  \frac{\dd x\,\dd z}{|v(H)-v(W)|}.
 \label{pin:chart-measure}
\end{equation}

\begin{lemma}[Resonance charts covering all target points]
\label{pin:all-charts}
For every $x_*\in\mathbb T_{2\pi}$, there is a nontrivial resonant
quartet whose first leg is $x_*$ and whose other three velocities are
pairwise distinct. Consequently the circle admits a finite cover by
target intervals satisfying
\eqref{pin:chart}--\eqref{pin:chart-measure}.
\end{lemma}

\begin{proof}
An explicit symmetric quartet treats every target except four points.
At those points some velocities coincide, so we construct an asymmetric
quartet instead. Reflection, translation, and compactness then give
the finite cover. First suppose that
\[
                  x_*\notin\mathcal E
                  :=\{0,\pi/2,\pi,3\pi/2\}.
\]
Take
\begin{equation}
                       (x,\pi-x,-x,\pi+x).
 \label{pin:symmetric-quartet}
\end{equation}
The four legs are distinct and conserve periodic momentum; energy
conservation follows from the evenness of $\omega$. Write
\begin{equation}
 a=\omega'(x)=\frac{d\sin x}{\omega(x)},\qquad
 b=\omega'(\pi-x)=\frac{d\sin x}{\omega(\pi-x)}.
 \label{pin:ab}
\end{equation}
The numbers $a,b$ are nonzero, have the same sign, and are unequal.
The non-target velocities $b,-a,-b$ are thus pairwise distinct.
In particular, at the center of this chart,
\begin{equation}
                   F_y=2b,\qquad
                   H_z=\frac{b-a}{2b},\qquad W_z=\frac{-b-a}{2b}.
 \label{pin:symmetric-derivatives}
\end{equation}

To treat the four exceptional points, temporarily allow an auxiliary
signed parameter $-1/2<\widetilde d<1/2$, $\widetilde d\ne0$, and set
\begin{equation}
 \begin{gathered}
 u=\sqrt{1-2\widetilde d},\qquad c_0=1-u,
 \qquad q=\sqrt{4+4c_0-7c_0^2},\\
 A=\frac{c_0+q}{2},\qquad B=\frac{q-c_0}{2}.
 \end{gathered}
 \label{pin:special-parameters}
\end{equation}
Then $0<u<\sqrt2$, $1-\sqrt2<c_0<1$, $c_0\ne0$, and
\begin{align}
 2\widetilde d&=c_0(2-c_0),&
 A-B&=c_0,& A+B&=q,\notag\\
 AB&=u(1+2c_0)>0,&
 A^2+B^2&=2+2c_0-3c_0^2.
 \label{pin:special-algebra-a}
\end{align}
Since $1+2c_0>3-2\sqrt2>0$, we have $q>|c_0|$ and hence $A,B>0$.
Moreover,
\begin{align}
 (1-A^2)^2+(1-B^2)^2&=c_0^2(2-c_0)^2=(2\widetilde d)^2,
 \label{pin:special-unit-circle}\\
 (2+c_0)^2-q^2&=8c_0^2>0,\notag\\
 q^2-(2-c_0)^2&=8c_0(1-c_0).
 \label{pin:special-signs}
\end{align}
Thus $B<1$, while $A>1$ if $\widetilde d>0$ and $A<1$ if
$\widetilde d<0$. Equation \eqref{pin:special-unit-circle} determines
a unique angle $y$ such that
\begin{equation}
              \cos y=\frac{1-A^2}{2\widetilde d},\qquad
              \sin y=\frac{1-B^2}{2\widetilde d}.
 \label{pin:special-y}
\end{equation}
For positive parameter, $y\in(\pi/2,\pi)$; for negative parameter,
$y\in(\pi,3\pi/2)$. With auxiliary dispersion
$\omega_{\widetilde d}(x)=\sqrt{1-2\widetilde d\cos x}$, we have
\[
 \omega_{\widetilde d}(0)=u,\quad
 \omega_{\widetilde d}(y)=A,\quad
 \omega_{\widetilde d}(\pi/2)=1,\quad
 \omega_{\widetilde d}(y-\pi/2)=B.
\]
The equality $u+A=1+B$ shows that
\begin{equation}
                          (0,y,\pi/2,y-\pi/2)
 \label{pin:special-quartet}
\end{equation}
conserves both periodic momentum and energy.

Put $w=y-\pi/2$. The four velocities are
$0,v_y,\widetilde d,v_w$, respectively, where
\begin{equation}
              v_y=\frac{1-B^2}{2A},\qquad
              v_w=\frac{A^2-1}{2B}.
 \label{pin:special-velocities}
\end{equation}
Factoring with \eqref{pin:special-algebra-a} gives
\begin{equation}
 \begin{aligned}
 v_y-v_w&=-\frac{c_0u(A+B)}{2AB},\\
 v_y-\widetilde d&=-\frac{c_0u(A+1)}{2A},\\
 v_w-\widetilde d&=\frac{c_0u(1-B)}{2B}.
 \end{aligned}
 \label{pin:velocity-differences}
\end{equation}
Every right-hand side is nonzero. Hence
\begin{equation}
 \begin{cases}
 0<v_y<\widetilde d<v_w,&\widetilde d>0,\\
 v_w<\widetilde d<0<v_y,&\widetilde d<0.
 \end{cases}
 \label{pin:velocity-order}
\end{equation}
In particular, all four velocities are pairwise distinct.

Return now to the fixed positive parameter $d$. With $\widetilde d=d$,
the quartet \eqref{pin:special-quartet} gives a chart at target $0$.
Exchanging the incoming and outgoing pairs writes it as
$(\pi/2,w,0,y)$ and gives a chart at target $\pi/2$. Reflecting all
four legs gives a chart at target $3\pi/2$. Finally, construct
\eqref{pin:special-quartet} with $\widetilde d=-d$ and translate all
four legs by $\pi$. The identities
\begin{equation}
          \omega_d(s+\pi)=\omega_{-d}(s),\qquad
          \omega_d'(s+\pi)=\omega_{-d}'(s)
 \label{pin:parameter-shift}
\end{equation}
place the resulting quartet on the resonance surface for the positive
parameter $d$ and give a chart at target $\pi$. Both momentum sums
increase by $2\pi$, so periodic momentum conservation is preserved.
This covers all four exceptional points. For each target, first choose
$x_*\in I'\Subset I$ inside its resonance chart $I\times J$, then
choose a finite subcover from these inner intervals $I'$. The averaging
kernels below are defined on the outer intervals $I$, while uniform
estimates are taken on $\overline{I'}$; subsequent chart restrictions
leave the chosen cover of the circle intact.
\end{proof}

In this atlas, each of the non-target maps $z,H(x,z),W(x,z)$ permits
a change of variable in $z$. For integrable $\varphi$, the resonance
average in \eqref{pin:averaged-invariance} can therefore be written as
the integral of $\varphi$ against a smooth kernel, with target
derivatives falling on that kernel. The resonance identity makes the
function equal to its average. To cover the finite measurable class of
Theorem~\ref{pin:classification}, the next proof first uses level sets
to establish boundedness, and only then integrates function values.

\begin{lemma}[From finite measurability to smoothness]
\label{pin:regularization}
Every measurable function that is finite almost everywhere and satisfies
\eqref{pin:invariance} belongs to $L^\infty(\mathbb T_{2\pi})$ and
has a smooth periodic representative. If $\varphi\in L^1$ is already
known, this representative satisfies, for each integer $r\ge0$,
\begin{equation}
                  \|\varphi\|_{C^r}\le C_{d,r}\|\varphi\|_{L^1(m)}.
 \label{pin:smoothing-bound}
\end{equation}
\end{lemma}

\begin{proof}
Fix inner and outer intervals and a rectangle from
Lemma~\ref{pin:all-charts}, denoted by $I'\Subset I$ and $I\times J$,
with $|H_z|\ge c_H>0$ and $|W_z|\ge c_W>0$. Set
$E_M=\{s:|\varphi(s)|>M\}$. Since the circle has finite measure and
$\varphi$ is finite almost everywhere, $|E_M|\to0$. For each fixed
$x\in I$, one-dimensional changes of variable give
\begin{align*}
 |\{z\in J:z\in E_M\}|&\le |E_M|,\\
 |\{z\in J:H(x,z)\in E_M\}|&\le c_H^{-1}|E_M|,\\
 |\{z\in J:W(x,z)\in E_M\}|&\le c_W^{-1}|E_M|.
\end{align*}
Choose a finite $M$ so that the sum of these bounds is less than
$|J|/2$. For each $x$, there is then a set of $z$ of measure at least
$|J|/2$ on which the three non-target values have absolute value at
most $M$. Integrability of the function has not been used. Apply Fubini
only to the indicator of the set where the invariant identity fails.
By \eqref{pin:chart-measure}, for almost every $x$ the identity holds
for almost every $z$. Thus, for almost every $x\in I$, the positive
measure set just constructed contains a $z$ at which the identity
holds, and
\[
       |\varphi(x)|
       =|\varphi(z)+\varphi(W(x,z))-\varphi(H(x,z))|\le3M.
\]
The same bound holds on $I'$. The finitely many inner intervals cover
the circle, so taking the maximum of their thresholds $M$ gives the
global $L^\infty$ bound.

We may now use $L^1$ regularization. Choose
$\chi\in C_c^\infty(J)$ with $\chi\ge0$ and
$\int\chi(z)\dd z=1$. The invariant relation gives, almost everywhere,
\begin{equation}
 \varphi(x)=\int\chi(z)\varphi(z)\dd z
   +\int\chi(z)\varphi(W(x,z))\dd z
   -\int\chi(z)\varphi(H(x,z))\dd z.
 \label{pin:averaged-invariance}
\end{equation}
Change variables along $W$ and $H$ in the last two terms. For example,
if $z=Z_H(x,s)$ is the inverse of $s=H(x,z)$, the corresponding
kernel is
\[
             K_H(x,s)=
             \frac{\chi(Z_H(x,s))}{|H_z(x,Z_H(x,s))|}.
\]
The support of $\chi$ has positive distance from $\partial J$, and
$|H_z|\ge c_H$, so the distance between $H(x,\supp\chi)$ and
$H(x,\partial J)$ is uniformly positive. The inverse function theorem
gives uniform bounds on all derivatives of $Z_H$ for
$x\in\overline{I'}$. Consequently $K_H$ vanishes near the boundary
of the image interval and extends smoothly by zero to a kernel on the
circle, with
\[
 \sup_{x\in\overline{I'},\,s\in\mathbb T_{2\pi}}
       |\partial_x^r K_H(x,s)|\le C_r\qquad(r\ge0).
\]
The $W$ term is treated in the same way; the first kernel is $\chi$,
independent of $x$. Thus the right-hand side of
\eqref{pin:averaged-invariance} is smooth, with each derivative bounded
by $C_r\|\varphi\|_1$. The finite cover gives
\eqref{pin:smoothing-bound}. Local smooth representatives agree almost
everywhere on overlaps and hence, by continuity, everywhere there.
They therefore define a unique global smooth periodic representative.

Lemma~\ref{pin:leg-ac} shows that replacing the function by this
representative preserves the invariant identity almost everywhere on
all of $\Sigma$. The collision difference is continuous on the regular
surface, and every regular open patch has positive measure. The
identity therefore holds pointwise at every regular resonant point.
\end{proof}

\subsection{Second- and third-order compatibility on the symmetric resonance branch}
\label{pin:jets}

We next prove that a $C^1$ periodic invariant must have the form
$A+B\omega$. Away from $\mathcal E$, the symmetric charts of
Lemma~\ref{pin:all-charts} and \eqref{pin:averaged-invariance} already
give local smoothness, so the second and third derivatives below require
only the initial $C^1$ assumption. If the identity is initially known
only almost everywhere on $\Sigma$, continuity first makes it pointwise
on these charts. We treat the odd and even parts separately. For the
odd part, two differentiations leave a nonzero coefficient times
$\varphi''$. For the even part, we subtract a multiple of the
third-order energy identity from the third-order invariant identity to
eliminate the unknown derivative of the resonance branch. The resulting
ordinary differential equation is solved on two open intervals, whose
solutions are then matched by endpoint regularity and periodicity.

Reflection of all four legs, $x_i\mapsto-x_i$, preserves resonance.
Hence the odd and even parts separately satisfy the invariant relation.
First let $\varphi$ be odd. On \eqref{pin:symmetric-quartet}, the
invariant relation gives
\begin{equation}
                  \varphi(\pi-x)=-\varphi(x),\qquad
                  \varphi(\pi+x)=\varphi(x).
 \label{pin:odd-relations}
\end{equation}
Fix $x\notin\mathcal E$ and take the following curve in its resonance
chart:
\begin{equation}
 z(t)=-x+t,\qquad y(t)=H(x,z(t)),\qquad
 w(t)=x+y(t)-z(t).
 \label{pin:transverse-curve}
\end{equation}
Write $y'(0)=\alpha$ and $y''(0)=\beta$. Then
$w'(0)=\alpha-1$ and $w''(0)=\beta$. By
\eqref{pin:odd-relations} and oddness, the derivatives at the four
base points satisfy
\begin{align*}
 \varphi'(\pi-x)&=\varphi'(-x)=\varphi'(\pi+x)=\varphi'(x),\\
 \varphi''(\pi-x)&=\varphi''(-x)=-\varphi''(x),\qquad
 \varphi''(\pi+x)=\varphi''(x).
\end{align*}
Differentiate
$\varphi(x)+\varphi(y(t))-\varphi(z(t))-\varphi(w(t))=0$
twice. The two $\beta\varphi'$ terms cancel, leaving
\begin{equation}
 0=[1-\alpha^2-(\alpha-1)^2]\varphi''(x)
   =\frac{b^2-a^2}{2b^2}\varphi''(x).
 \label{pin:odd-second}
\end{equation}
Since $a^2\ne b^2$ in \eqref{pin:ab}, we obtain $\varphi''=0$ on
$\mathbb T_{2\pi}\setminus\mathcal E$. The $C^1$ assumption matches
the constant values of $\varphi'$ on the intervening intervals.
Periodicity forces their common slope to vanish, and oddness then
forces the function itself to vanish.

Now let $\varphi$ be even and fix
$x\in(0,\pi)\setminus\{\pi/2\}$. Along
\eqref{pin:transverse-curve}, write
\[
 y'(0)=\alpha,\quad y''(0)=\beta,\quad y'''(0)=\gamma,
 \qquad c_a=\omega''(x),\quad c_b=\omega''(\pi-x),\quad r=a/b.
\]
The first- and second-order energy identities give, respectively,
\begin{equation}
                    \alpha=\frac{1-r}{2},\qquad
                    \beta=\frac{c_a+r c_b}{2b}.
 \label{pin:branch-jets}
\end{equation}
Define $g(x)=\varphi'(x)/\omega'(x)$. The first-order invariant
identity is
\[
                   0=(2\alpha-1)\varphi'(\pi-x)+\varphi'(x),
\]
and therefore
\begin{equation}
       g(\pi-x)=g(x),\qquad
       g'(\pi-x)=-g'(x),\qquad g''(\pi-x)=g''(x).
 \label{pin:g-reflection}
\end{equation}
The third-order energy and invariant identities are
\begin{align}
 0={}&2b\gamma+3c_b\beta
   +\omega'''(\pi-x)[\alpha^3+(\alpha-1)^3]+\omega'''(x),
 \label{pin:energy-third}\\
 0={}&2\varphi'(\pi-x)\gamma+3\varphi''(\pi-x)\beta
   +\varphi'''(\pi-x)[\alpha^3+(\alpha-1)^3]+\varphi'''(x).
 \label{pin:invariant-third}
\end{align}
Subtract $g(x)$ times \eqref{pin:energy-third} from
\eqref{pin:invariant-third}. This eliminates the third branch
derivative $\gamma$. Using $\varphi'=\omega'g$ and
\eqref{pin:g-reflection}, we have
\begin{align}
 \varphi''(\pi-x)-gc_b&=-bg',\notag\\
 \varphi'''(\pi-x)-g\omega'''(\pi-x)&=-2c_b g'+bg'',\notag\\
 \varphi'''(x)-g\omega'''(x)&=2c_a g'+ag''.
 \label{pin:third-subtraction}
\end{align}
With $H_3=\alpha^3+(\alpha-1)^3=-r(3+r^2)/4$, the remaining
identity reads
\[
             -3b\beta g'+(-2c_bg'+bg'')H_3+2c_ag'+ag''=0.
\]
Substitution of \eqref{pin:branch-jets} and simplification yield
\begin{equation}
                  a(1-r^2)g''+2(c_a+r^3c_b)g'=0.
 \label{pin:g-ode}
\end{equation}
Since $b'(x)=-c_b$, direct differentiation gives
\begin{equation}
 \frac{\dd}{\dd x}\log\left|\frac{a^2b^2}{b^2-a^2}\right|
        =\frac{2(c_a+r^3c_b)}{a(1-r^2)}.
 \label{pin:integrating-factor}
\end{equation}
Thus \eqref{pin:g-ode} is equivalent to
\begin{equation}
                 \left(\frac{a^2b^2}{b^2-a^2}g'\right)'=0.
 \label{pin:conservative-ode}
\end{equation}
For the dispersion \eqref{pin:dispersion}, the integrating factor
simplifies explicitly:
\begin{equation}
 a^2=\frac{d^2\sin^2x}{1-2d\cos x},\quad
 b^2=\frac{d^2\sin^2x}{1+2d\cos x},\quad
 \frac{a^2b^2}{b^2-a^2}=-\frac{d\sin^2x}{4\cos x}.
 \label{pin:explicit-factor}
\end{equation}
On each of the intervals $(0,\pi/2)$ and $(\pi/2,\pi)$, we obtain
\begin{equation}
                   g(x)=A_0+\frac{B_0}{\sin x},\qquad
                   \varphi'(x)=\frac{d(A_0\sin x+B_0)}{\omega(x)}.
 \label{pin:local-ode-solutions}
\end{equation}
For now the constants may differ between the two intervals. An even
periodic $C^1$ function satisfies $\varphi'(0)=\varphi'(\pi)=0$.
Since $\omega(0),\omega(\pi)>0$, the two endpoints force
$B_0=0$ on their respective intervals. At $\pi/2$ we have
$\omega'(\pi/2)=d\ne0$, so continuity of $\varphi'$ also matches
the two values of $A_0$. Continuity of $\varphi$ matches the constants
of integration. Hence $\varphi=A+B\omega$ on $[0,\pi]$, and evenness
and periodicity extend the identity to the circle.

This proves the $C^1$ classification. Together with
Lemma~\ref{pin:regularization}, it proves
Theorem~\ref{pin:classification}. Since $\omega(0)\ne\omega(\pi)$,
the two coefficients are unique. The converse follows directly from
\eqref{pin:resonance}.

\begin{remark}[The interval equation and endpoint resonances]
\label{pin:endpoint-test}
Eliminating the $B_0$ term in \eqref{pin:local-ode-solutions} requires
the endpoint conditions. For example, the continuous even periodic
function
\[
 \Phi(x)=d\int_0^{\dist(x,2\pi\mathbb Z)}\frac{\dd s}{\omega(s)}
\]
gives $g=1/\sin x$ on $(0,\pi)$ and satisfies
\eqref{pin:g-reflection} and \eqref{pin:g-ode}, but has corners at
$0,\pi$. For the positive-parameter value of $y$ in
\eqref{pin:special-quartet}, put $w=y-\pi/2\in(0,\pi/2)$.
Its collision difference is
\[
 \Delta\Phi
   =d\int_0^w\left(\frac1{\omega(s+\pi/2)}
                         -\frac1{\omega(s)}\right)\dd s<0.
\]
Here $\omega(s+\pi/2)>\omega(s)$. Thus the endpoint regularity
provided by charts at every target excludes this branch, which the
third-order equation on the open intervals alone would retain.
\end{remark}

\begin{corollary}[Positive affine-frequency denominators]
\label{pin:rj}
Suppose $0<W<\infty$ almost everywhere and $1/W$ satisfies
\eqref{pin:invariance}. Then $W=(A+B\omega)^{-1}$ almost everywhere,
with $A,B\in\R$. If also $W\in L^1(m)$, then
\begin{equation}
                         \min_x(A+B\omega(x))>0.
 \label{pin:rj-positive}
\end{equation}
\end{corollary}

\begin{proof}
Theorem~\ref{pin:classification} applies to the finite measurable
function $1/W$. Since $\omega$ is nonconstant, the identity
$\operatorname{Im}A+\operatorname{Im}B\,\omega=0$ forces $A,B$ to
be real. The continuous denominator is positive almost everywhere and
therefore nonnegative everywhere. If it vanishes somewhere, a nonzero
affine function of frequency can vanish only at a maximum or minimum
of $\omega$. Both extrema are nondegenerate and quadratic, so the
reciprocal has a nonintegrable quadratic pole. An identically zero
denominator contradicts the almost-everywhere finiteness and positivity
of $W$. This proves \eqref{pin:rj-positive}.
\end{proof}

The corollary assumes the resonance invariant identity for $1/W$.
Starting from a nonlinear stationary state, one must still justify the
collision integral and entropy production identity in the relevant
function class to obtain that assumption. Condition
\eqref{pin:rj-positive} allows $B$ to be positive, zero, or negative.

\subsection{Resonance averaging and a positive spectral gap}
\label{pin:gap}

Let $a(x_0,x_1,x_2,x_3)$ be a smooth positive weight on the compact
four-dimensional torus such that
\begin{equation}
                         0<a_-\le a\le a_+<\infty.
 \label{pin:weight}
\end{equation}
In the reduced coordinates it is always evaluated at
$a(x,y,z,x+y-z)$. The lower bound in this subsection requires no
additional symmetry. When relating it to the linearized collision
operator, we will require invariance of $a$ under exchange of the two
incoming legs, exchange of the two outgoing legs, and exchange of the
two pairs. Define the extended quadratic form
\begin{equation}
 \mathcal D_a(\varphi)
       =\int_\Sigma a|\Delta\varphi|^2\dd\Xi\in[0,\infty],
 \qquad \varphi\in L^2(m).
 \label{pin:extended-form}
\end{equation}
Lemma~\ref{pin:leg-ac} makes this well-defined on $L^2$ equivalence
classes. Let $P_0$ be the orthogonal projection in $L^2(m)$ onto
$\Span\{1,\omega\}$.

\begin{theorem}[Positive spectral gap]
\label{pin:gap-theorem}
There is $\lambda_{d,a}>0$ such that every $\varphi\in L^2(m)$ satisfies
\begin{equation}
             \mathcal D_a(\varphi)
               \ge\lambda_{d,a}\|(I-P_0)\varphi\|_{L^2(m)}^2.
 \label{pin:gap-inequality}
\end{equation}
The left-hand side may be infinite.
\end{theorem}

\begin{proof}
We construct a smooth averaging operator $K$ for which the dissipation
controls $\varphi-K\varphi$. A bounded sequence with vanishing
dissipation is then asymptotic to the image of a compact operator and
has a strongly convergent subsequence. We pass zero dissipation to its
limit and apply the classification of invariants.
Choose a nonnegative smooth partition of unity $\rho_i$ subordinate
to the finite inner-interval cover of Lemma~\ref{pin:all-charts}, with
$\supp\rho_i\Subset I'_i\Subset I_i$. Choose also
$\chi_i\in C_c^\infty(J_i)$, $\chi_i\ge0$, $\int\chi_i=1$.
For an arbitrary $\varphi\in L^2$, define
\begin{align}
 K_i\varphi(x)&=\int\chi_i(z)
  [\varphi(z)+\varphi(W_i(x,z))-\varphi(H_i(x,z))]\dd z,\notag\\
 K\varphi&=\sum_i\rho_iK_i\varphi.
 \label{pin:compact-average}
\end{align}
The change-of-variable argument in Lemma~\ref{pin:regularization}
shows that $K$ has a global smooth kernel and that, for each integer
$r\ge0$,
\begin{equation}
                     \|K\varphi\|_{C^r}\le C_{d,r}\|\varphi\|_2.
 \label{pin:compact-smoothing}
\end{equation}
In particular, $K:L^2\to L^2$ is compact. It need not be a projection
or a self-adjoint operator.

By definition and $\int\chi_i=1$,
\[
 \varphi(x)-K_i\varphi(x)
  =\int\chi_i(z)\Delta\varphi(x,H_i(x,z),z)\dd z.
\]
Applying Jensen's inequality twice gives
\begin{align}
 \|\varphi-K\varphi\|_2^2
 &\le\sum_i\int\rho_i(x)\chi_i(z)
       |\Delta\varphi(x,H_i(x,z),z)|^2m(\dd x)\dd z\notag\\
 &\le C_{d,a}\mathcal D_a(\varphi).
 \label{pin:defect-bound}
\end{align}
The last step uses \eqref{pin:chart-measure}: the Jacobian factors and
weights have uniform bounds on the compact supports, and finite
overlap contributes only a finite constant. For a general $L^2$
function, the three compositions on a chart are locally in $L^2$, so
Fubini justifies this step. The inequality is automatic when the
right-hand side is infinite.

If \eqref{pin:gap-inequality} failed, there would be
$\varphi_n\perp\Span\{1,\omega\}$ with $\|\varphi_n\|_2=1$ and
$\mathcal D_a(\varphi_n)\to0$. By \eqref{pin:defect-bound},
$\varphi_n-K\varphi_n\to0$ in $L^2$. Compactness gives a subsequence
$\varphi_n\to\varphi$ strongly in $L^2$, with $\|\varphi\|_2=1$ and
$P_0\varphi=0$. Pass to a further subsequence converging almost
everywhere in the one-dimensional variable. Lemma~\ref{pin:leg-ac}
transfers convergence simultaneously along all four legs to all of
$\Sigma$. Fatou's lemma gives
\[
                    \mathcal D_a(\varphi)
                    \le\liminf_n\mathcal D_a(\varphi_n)=0.
\]
Theorem~\ref{pin:classification} forces
$\varphi\in\Span\{1,\omega\}$, contradicting its unit norm and
orthogonality. This limiting step uses the coarea measure on all regular
resonances; the finite averaging charts are used only for the compactness
estimate \eqref{pin:defect-bound}.
Finally, $\mathcal D_a\ge a_-\mathcal D_1$ permits the choice
$\lambda_{d,a}=a_-\lambda_{d,1}$, which also yields a uniform lower
bound as the weight varies over a compact positive parameter family.
\end{proof}

\subsection{Critical resonances and cancellation in the collision difference}
\label{pin:critical}

Resonance averaging has already given the positive spectral gap. To
construct a dense domain for \eqref{pin:extended-form}, we must also
prove that smooth test functions have finite energy near critical
points of the energy. The key is the common factor in
\eqref{pin:double-difference} below:
$F=-uvG_\omega$ and $\Delta\varphi=-uvG_\varphi$.
Forming the collision difference before taking the regularized energy
$\delta$ limit preserves this cancellation within the same integral.
The factorization uses second-order calculus for $\varphi\in C^2$;
general $L^2$ functions will subsequently be treated through the
maximal closed difference operator. We first show that every critical
resonance admits the required coordinates. Direct calculation gives
\begin{equation}
 \omega'(x)=\frac{d\sin x}{\omega(x)},\qquad
 \omega''(x)=\frac{d[\cos x-d(1+\cos^2x)]}{\omega(x)^3}.
 \label{pin:curvature}
\end{equation}
The polynomial $h(t)=t-d(1+t^2)$ is strictly increasing on $[-1,1]$,
since $h'(t)=1-2dt>0$, and $h(-1)<0<h(1)$. Its unique zero
$t_*\in(0,1)$ is simple. Thus $\omega'$ first increases strictly
and then decreases strictly on $(0,\pi)$, and has a corresponding
unique minimum on $(\pi,2\pi)$. Both zeros of the curvature are simple;
for example, $x_* =\arccos t_*$ satisfies
\begin{equation}
             \omega'''(x_*)=
             -\frac{d(1-2d\cos x_*)\sin x_*}{\omega(x_*)^3}\ne0.
 \label{pin:simple-inflection}
\end{equation}

\begin{lemma}[Critical points and their local coordinates]
\label{pin:critical-charts}
Every critical quartet in the resonance set $\{F=0\}$ is a trivial
permutation. Near each critical point, after permuting the legs and
taking local real lifts, one can choose coordinates
\begin{equation}
       x=z+u,\qquad y=z+v,\qquad w=z+u+v
 \label{pin:critical-lift}
\end{equation}
in which, for every $C^2$ function $\varphi$,
\begin{equation}
 F=-uvG_\omega,\qquad \Delta\varphi=-uvG_\varphi,
 \qquad
 G_\varphi(z,u,v)=\int_0^1\int_0^1
                \varphi''(z+su+tv)\dd s\dd t.
 \label{pin:double-difference}
\end{equation}
The neighborhood can be chosen to have one of the following two forms:
\begin{enumerate}
 \item $|G_\omega|$ is uniformly bounded away from zero, and the
       center has $u=v=0$;
 \item $|\partial_zG_\omega|$ is uniformly bounded away from zero,
       so that $q=G_\omega$ can be used as a coordinate.
\end{enumerate}
\end{lemma}

\begin{proof}
The unimodality following from \eqref{pin:curvature} shows that each
velocity value has at most two preimages on the circle. If the preimages
are distinct, they have curvatures of opposite signs and distinct
frequencies: positive and negative velocities lie in their respective
semicircles, on each of which $\omega$ is strictly monotone; zero
velocity occurs only at $0,\pi$. On the other hand,
\begin{equation}
       \nabla F=0\quad\Longleftrightarrow\quad
       \omega'(x)=\omega'(y)=\omega'(z)=\omega'(w).
 \label{pin:critical-velocities}
\end{equation}
The four legs of a critical quartet thus take at most two values $p,q$.
If $p$ occurs $j,k$ times in the incoming and outgoing pairs,
respectively, energy conservation gives
\[
                    (j-k)[\omega(p)-\omega(q)]=0.
\]
Distinct points have distinct frequencies, so $j=k$. The two multisets
agree, and the quartet is a trivial permutation. The case of four
identical legs is included.

For the trivial pairing $x=z,y=w$, choose
\eqref{pin:critical-lift}, with $u=0$ at the center. If the other two
points are distinct, choose a nonzero local real representative
$v_0\ne0$; if all four legs coincide, take $u_0=v_0=0$. The change
of coordinates has determinant of absolute value $1$, and periodic
functions are represented by their periodic real lifts. Two applications
of the fundamental theorem of calculus give
\eqref{pin:double-difference}. Exchanging the two incoming legs reduces
the other trivial pairing to this case.

At a critical point with $u=0,v\ne0$,
\begin{align}
 G_\omega(z,0,v)&=\frac{\omega'(z+v)-\omega'(z)}{v}=0,\notag\\
 \partial_zG_\omega(z,0,v)
       &=\frac{\omega''(z+v)-\omega''(z)}{v}\ne0.
 \label{pin:separated-critical}
\end{align}
The last inequality follows because the two preimages of the same
velocity have curvatures of opposite signs. At $u=v=0$,
$G_\omega=\omega''(z)$. If this vanishes, then
$\partial_zG_\omega=\omega'''(z)\ne0$ by
\eqref{pin:simple-inflection}, giving the second type of chart. If it
does not vanish, we obtain the first type. Every quartet with four
identical legs is critical, so the first type must also be included.
Compactness of the critical set gives a finite cover by such
neighborhoods.
\end{proof}

Equation \eqref{pin:double-difference} is a double-difference identity.
For example, taking $\varphi(t)=t^2$ only in a local real lift gives
\[
 (z+u)^2+(z+v)^2-z^2-(z+u+v)^2=-2uv,
 \qquad G_\varphi=2.
\]
This local, nonperiodic calculation explains the factor $uv$: the four
terms cancel the constant and linear parts, leaving only the mixed
second-order term. Averaging the second derivative in
\eqref{pin:double-difference} gives the same structure for any $C^2$
function.

\begin{proposition}[Finite energy of smooth functions and the regularized limit]
\label{pin:mollifier}
For $\varphi\in C^2(\mathbb T_{2\pi})$,
\begin{equation}
 \mathcal D_a(\varphi)\le C_{d,a}\|\varphi\|_{C^2}^2,
 \qquad
 \int_\Sigma a|\Delta\varphi|\dd\Xi
       \le C_{d,a}\|\varphi\|_{C^2}.
 \label{pin:finite-energy}
\end{equation}
Let $\rho_\eta\in C_c^\infty(\R)$ be any family of nonnegative
functions satisfying
\begin{equation}
       \int\rho_\eta(s)\dd s=1,\qquad
       \supp\rho_\eta\subset[-\eta,\eta].
 \label{pin:mollifier-assumptions}
\end{equation}
Then, for $\varphi,\psi\in C^2$,
\begin{equation}
 \begin{split}
 &\lim_{\eta\downarrow0}
 \int_{\mathbb T_{2\pi}^3}
      a\Delta\varphi\,\overline{\Delta\psi}\,
      \rho_\eta(F)\,m(\dd x)m(\dd y)m(\dd z)\\
 &\hspace{25mm}=\int_\Sigma a\Delta\varphi\,
                               \overline{\Delta\psi}\dd\Xi.
 \end{split}
 \label{pin:mollified-form}
\end{equation}
No evenness, fixed profile, or uniform $O(\eta^{-1})$ height bound is
required of $\rho_\eta$.
\end{proposition}

\begin{proof}
We first estimate the collision difference in the two types of critical
chart and then treat the regular region away from the critical set.
In the second type, the nontrivial resonance branch is $G_\omega=0$.
The first type has only trivial zero branches, and its regularized
integral must separately be shown to vanish.
In a chart of the second type, change variables to $q=G_\omega$.
On the nontrivial regular branch $q=0$, $uv\ne0$,
\eqref{pin:double-difference} gives
\begin{align}
 |\Delta\varphi|^2\dd\Xi
   &=(2\pi)^{-3}
      \frac{|uv|\,|G_\varphi|^2}{|\partial_zG_\omega|}\dd u\dd v,
 \label{pin:q-form-density}\\
 |\Delta\varphi|\dd\Xi
   &=(2\pi)^{-3}
      \frac{|G_\varphi|}{|\partial_zG_\omega|}\dd u\dd v.
 \label{pin:q-linear-density}
\end{align}
Both expressions are locally integrable, since
$|G_\varphi|\le\|\varphi''\|_\infty$. The collision difference
vanishes on trivial regular branches. Charts of the first type contain
only trivial zero branches and hence contribute nothing to the form.
Regular compact sets away from the critical set have finite coarea
measure, so a finite cover proves \eqref{pin:finite-energy}.

Choose a smooth partition of unity subordinate to these charts and the
remaining regular region. In a chart of the second type, the
regularized integral at fixed $u,v$ has the form
\begin{equation}
             u^2v^2\int A(q,u,v)\rho_\eta(-uvq)\dd q,
 \label{pin:q-inner}
\end{equation}
where $A$ includes the weight, cutoff,
$G_\varphi\overline{G_\psi}$, and Jacobian factor; it is continuous
and uniformly bounded. The absolute value of \eqref{pin:q-inner} is
bounded by $C|uv|$, using
$\int\rho_\eta(-uvq)\dd q\le |uv|^{-1}$. For each fixed
$uv\ne0$, the one-dimensional approximate identity gives
$|uv|A(0,u,v)$. Dominated convergence on the bounded $u,v$ region
gives the sesquilinear version of \eqref{pin:q-form-density}.
The axes $uv=0$ form a two-dimensional null set, and the uniform
$C|uv|$ bound also controls their neighborhoods.

A chart of the first type is centered at $u=v=0$. After shrinking
the neighborhood, set $V=vG_\omega(z,u,v)$; its derivative with
respect to $v$ is bounded away from zero. Then
\begin{equation}
      F=-uV,\qquad
      \Delta\varphi=-uV\frac{G_\varphi}{G_\omega},\qquad
      |\Delta\varphi\Delta\psi|\le C u^2V^2.
 \label{pin:nonzero-G}
\end{equation}
Take a rectangle $[-L,L]^2$ containing the projection of this chart and set
\[
 J_j(\eta)=\int_{[-L,L]^2}|uV|^j\rho_\eta(-uV)\dd u\dd V.
\]
For fixed $u\ne0$, a one-dimensional change of variables bounds the
slice of $J_1$ by $L$; the support condition gives the additional bound
$\eta/|u|$. Splitting the integral into $|u|\le\delta$ and
$|u|>\delta$ yields $J_1(\eta)\le2L\delta+2L\eta/\delta$.
Choosing $\delta=\sqrt\eta$, we obtain
$J_1(\eta)\le4L\sqrt\eta$. On the support of the kernel,
$|uV|\le\eta$, so $J_2(\eta)\le\eta J_1(\eta)$. The coordinate
Jacobian is bounded, and the contribution of this chart satisfies
\begin{align}
 \left|\text{regularized integral}\right|
 &\le C J_2(\eta)\le C\eta^{3/2}\longrightarrow0.
 \label{pin:product-mollifier}
\end{align}
Only \eqref{pin:mollifier-assumptions} was used, so the argument applies
to any family stated in the proposition. In the remaining regular
region, the usual coarea change of variables applies; compact sets away
from $F=0$ eventually contribute nothing. Summing finitely many pieces
proves \eqref{pin:mollified-form}.
\end{proof}

This limit regularizes the critical energy after the collision
difference has been formed. It does not assign a finite measure limit
to $\rho_\eta(F)$ alone, without cancellation. For smooth collision
differences there is no residual critical contribution to the form.

\subsection{The maximal closed form, self-adjoint realization, and linearization normalization}
\label{pin:closed-realization}

Set
\begin{equation}
 \mathcal H=L^2(\mathbb T_{2\pi},m),\qquad
 \mathcal G_a=L^2(\Sigma,a\dd\Xi),\qquad
 \Dom R_a=\{\varphi\in\mathcal H:\Delta\varphi\in\mathcal G_a\},
 \quad R_a\varphi=\Delta\varphi.
 \label{pin:maximal-difference}
\end{equation}
The domain contains every $L^2$ function whose collision difference is
square integrable; this is the maximal difference operator. The four
individual pullbacks need not belong to $\mathcal G_a$, since the
domain allows cancellation between them. Lemma~\ref{pin:leg-ac}
ensures that the operator is independent of representatives.

\begin{theorem}[Closed form and spectral-gap realization]
\label{pin:closed-theorem}
The operator $R_a$ is closed and densely defined, with
$C^2(\mathbb T_{2\pi})\subset\Dom R_a$. For every $\gamma>0$, the form
\begin{equation}
 \mathfrak q_a(\varphi,\psi)=\frac\gamma4
                 \langle R_a\varphi,R_a\psi\rangle_{\mathcal G_a},
 \qquad \Dom\mathfrak q_a=\Dom R_a
 \label{pin:closed-form}
\end{equation}
is closed and densely defined. Its associated nonnegative self-adjoint
operator is
\begin{equation}
 L_a=\frac\gamma4R_a^*R_a,\qquad
 \Dom L_a=\{\varphi\in\Dom R_a:R_a\varphi\in\Dom R_a^*\}.
 \label{pin:selfadjoint-operator}
\end{equation}
Moreover,
\begin{equation}
 \Ker L_a=\Span\{1,\omega\},\qquad
 \mathfrak q_a(\varphi,\varphi)
   \ge\frac\gamma4\lambda_{d,a}\|(I-P_0)\varphi\|_2^2.
 \label{pin:operator-gap}
\end{equation}
\end{theorem}

\begin{proof}
Proposition~\ref{pin:mollifier} gives $C^2\subset\Dom R_a$, hence
density of the domain. Suppose $\varphi_n\to\varphi$ in $\mathcal H$
and $R_a\varphi_n\to Z$ in $\mathcal G_a$. First take a subsequence
such that $\varphi_n\to\varphi$ almost everywhere in the
one-dimensional variable. Lemma~\ref{pin:leg-ac} gives simultaneous
almost-everywhere convergence of its four pullbacks on all of
$\Sigma$, so $R_a\varphi_n\to\Delta\varphi$ almost everywhere.
Take a further almost-everywhere convergent subsequence from the
$\mathcal G_a$-convergent sequence to obtain $Z=\Delta\varphi$.
Thus $\varphi\in\Dom R_a$, proving that the operator is closed.

The square of the graph norm is
$\|\varphi\|_2^2+(\gamma/4)\|R_a\varphi\|_{\mathcal G_a}^2$.
Closedness makes the domain a Hilbert space in this norm. For every
$f\in\mathcal H$, the Riesz representation theorem gives a unique
$u\in\Dom R_a$ satisfying
\[
      \langle u,v\rangle+\mathfrak q_a(u,v)=\langle f,v\rangle
                    \quad(v\in\Dom R_a).
\]
By the definition of the adjoint, this is equivalent to
$u\in\Dom L_a$ and $(I+L_a)u=f$. To verify density of the operator
domain as well, denote the solution operator by $Sf=u$. Taking $v=u$
gives $\|Sf\|_2\le\|f\|_2$; testing the solutions for two
right-hand sides against one another gives
\[
 \langle Sf,g\rangle
   =\langle Sf,Sg\rangle+\mathfrak q_a(Sf,Sg)
   =\langle f,Sg\rangle.
\]
Thus $S$ is bounded and self-adjoint. If $Sf=0$, the variational
identity makes $f$ orthogonal to the dense domain $\Dom R_a$, hence
$f=0$. Therefore $\overline{\Ran S}=(\Ker S)^{\perp}=\mathcal H$,
while $\Ran S=\Dom L_a$. We now know that $L_a$ is densely defined,
symmetric, and nonnegative, and that $I+L_a$ is surjective. For
$h\in\Dom L_a^*$, take $u=S(h+L_a^*h)$. Then
$(I+L_a^*)(h-u)=0$; this kernel is orthogonal to
$\Ran(I+L_a)=\mathcal H$, so $h=u\in\Dom L_a$. This proves
self-adjointness and identifies the operator associated with
\eqref{pin:closed-form}. Its kernel is $\Ker R_a$.
Theorem~\ref{pin:classification} identifies that kernel, and
Theorem~\ref{pin:gap-theorem} gives \eqref{pin:operator-gap}.
\end{proof}

The realization in Theorem~\ref{pin:closed-theorem} is determined by
the maximal integral domain \eqref{pin:maximal-difference}.
Proposition~\ref{pin:mollifier} shows that $C^2$ is a dense subspace
of finite-energy functions, where density refers to the ambient
$L^2$ space. Establishing that $C^2$ is a form core or an operator
core would require the corresponding approximation in graph norm.
In particular, $C^2\subset\Dom R_a$ does not by itself imply
$C^2\subset\Dom L_a$.

We conclude by relating this realization to the weak collision
expression. Suppose $a$ has the four-leg symmetries stated above.
For $\varphi,\psi\in C^2$, \eqref{pin:finite-energy} ensures that
each integral
\[
                    \int_\Sigma a|\Delta\varphi|\,|\psi(x_i)|\dd\Xi
\]
is finite. Exchanging the incoming legs leaves the collision difference
unchanged, whereas exchanging the incoming and outgoing pairs reverses
its sign. These changes in the reduced coordinates have determinant
of absolute value $1$ and preserve the coarea measure. We may therefore
expand the collision difference of the test function term by term to
obtain
\begin{equation}
       \mathfrak q_a(\varphi,\psi)
        =\gamma\int_\Sigma
                    a\Delta\varphi\,\overline{\psi(x)}\dd\Xi.
 \label{pin:weak-output}
\end{equation}
The same critical coordinates prove the regularized limit with a
single collision difference. In a $q=G_\omega$ chart,
$uv\int A(q,u,v)\rho_\eta(-uvq)\dd q$ is bounded in absolute value
by a constant. For each fixed $uv\ne0$, the inner integral has its
one-dimensional approximate-identity limit. The constant bound is
integrable on the bounded $u,v$ region, and the axes remain a
two-dimensional null set. In a $G_\omega\ne0$ chart, the contribution
is at most $C\sqrt\eta$. Consequently,
\begin{equation}
 \begin{split}
 &\lim_{\eta\downarrow0}\gamma
 \int_{\mathbb T_{2\pi}^3}
       a\Delta\varphi\,\overline{\psi(x)}\rho_\eta(F)\,\dd m^3\\
 &\hspace{25mm}=\gamma\int_\Sigma
                    a\Delta\varphi\,\overline{\psi(x)}\dd\Xi.
 \end{split}
 \label{pin:weak-mollifier}
\end{equation}
This proves the limit directly even for a noneven $\rho_\eta$;
symmetrization by exchanging the two pairs is not used at finite
$\eta$.

For fixed $\varphi\in C^2$, the pushforward of the finite complex
measure $\gamma a\Delta\varphi\dd\Xi$ along the first leg is
absolutely continuous with respect to $m$ by Lemma~\ref{pin:leg-ac}.
Thus \eqref{pin:weak-output} defines an $L^1(m)$ weak output. On
$C^2\cap\Dom L_a$, it agrees with the $L^2$ output of
\eqref{pin:selfadjoint-operator}.

The corresponding normalization in
\cite[equations (3.21) and (4.1)]{ALS2006} is
\begin{equation}
             a=\prod_{i=0}^3\omega(x_i)^{-2},\qquad
             \gamma=\frac{9\pi}{4},\qquad
 \mathfrak q_{\rm ALS}(\varphi,\varphi)
     =\frac{9\pi}{16}\int_\Sigma
             \prod_{i=0}^3\omega(x_i)^{-2}|\Delta\varphi|^2\dd\Xi.
 \label{pin:als-normalization}
\end{equation}
The angular normalization $m(\dd x)=\dd x/(2\pi)$ and the factor
$(2\pi)^{-3}$ after momentum elimination are already included in
\eqref{pin:measure}.

More generally, let $N=(\alpha+\beta\omega)^{-1}>0$ and set
\begin{equation}
                         a_N=\frac{\prod_{i=0}^3N(x_i)}
                                          {\prod_{i=0}^3\omega(x_i)}.
 \label{pin:physical-weight}
\end{equation}
The form in the physical relative perturbation $W=N+Ng$ is
\begin{equation}
 \mathfrak q_N(g,h)=\frac\gamma4
     \int_\Sigma a_N\Delta(g/N)\,\overline{\Delta(h/N)}\dd\Xi.
 \label{pin:physical-form}
\end{equation}
The positive bounded invertible multiplication $g\mapsto g/N$ makes
this form closed and densely defined, with kernel
$N\Span\{1,\omega\}$. The inequality
\[
 \dist_{L^2}(g,N\Span\{1,\omega\})
       \le\|N\|_\infty
          \dist_{L^2}(g/N,\Span\{1,\omega\})
\]
gives the corresponding positive spectral gap. If $(\alpha,\beta)$
ranges over a fixed compact positive parameter set, the lower bound
for the weight and the multiplication bounds are uniform, so the gap
constant can also be chosen uniformly.

To check the sign, write the factor in the weak four-wave collision expression as
\[
 \gamma(\prod\omega_i^{-1})\prod W_i\,\Delta(1/W).
\]
At $W=N+N^2\varphi$, the identity $\Delta(1/N)=0$ eliminates the
first variations of the four multiplicative weights, while
$\delta(1/W)=-\varphi$. The linear term is therefore
$-\gamma\int a_N\Delta\varphi$. The positive sign of the dissipation
operator agrees with \eqref{pin:weak-output}. The conversion between
$g$ and $\varphi$ is specified by \eqref{pin:physical-form}; the
positive multiplication is not a unitary transformation. In particular,
for $N=(\beta\omega)^{-1}$ one has
$a_N=\beta^{-4}\prod_i\omega_i^{-2}$, so in the $N^2\varphi$
variable the linear term is $-\beta^{-4}L_{\rm ALS}\varphi$, in
agreement with the temperature factor in
\cite[equation (3.20)]{ALS2006}.
The conclusions of this section concern collision invariants, the
dissipation gap, and its closed-form realization. A nonlinear diffusion
limit with spatial transport requires additional dynamical estimates.
